\section{Closure Density and Geometry}
\label{sec:closure_density}

\subsection{The Quantity Called Closure Density}

Section~\ref{sec:gravity_distance} introduced $L_{ab}$ as the additional
contribution to the path cost:
\begin{equation}
  \boxed{L_{ab} = \sum_k \phi_{abk}.}
\end{equation}
This represents the total contribution of triangular closures containing
the edge $(a,b)$.
Each triangle contributes to structural congestion,
and the sum measures how densely that edge is embedded in closure structure.
Intuitively, regions with many triangles are structurally dense;
regions with few triangles are structurally sparse.
\begin{equation}
  \boxed{\text{Closure density = density of structure.}}
\end{equation}
This interpretation follows directly from Vol.~2,
where triangular closure was identified as the minimal persistent structure.

\subsection{Why It Distorts Distance}

The modified edge cost was $w_{ab} = -\log S_{ab} + \epsilon L_{ab}$.
The first term favors highly permeable paths;
the second penalizes structurally congested regions.
Thus, large closure density increases traversal cost.
\begin{equation}
  \boxed{\text{Closure density effectively stretches distance.}}
\end{equation}
This reproduces the qualitative behavior associated with gravity:
dense regions become harder to traverse,
propagation paths bend around them,
and effective geometry is distorted.
Unlike general relativity, however, curvature is not introduced as a fundamental
geometric postulate. It emerges from the modification of path selection.

\subsection{Tensorization}

Closure density is directional.
An edge contributes differently depending on orientation.
To describe this, define a coarse-grained tensor over a region $U$:
\begin{equation}
  L_{\mu\nu}(x)
  = \left\langle L_{ab}\,\hat{e}^{(ab)}_\mu \hat{e}^{(ab)}_\nu \right\rangle_U,
\end{equation}
where $\hat{e}^{(ab)}_\mu$ is the direction vector determined internally
from the distance structure. No external embedding space is assumed.
\begin{equation}
  \boxed{L_{\mu\nu} = \text{closure tensor.}}
\end{equation}
This tensor encodes directional structural congestion and therefore determines
anisotropic modifications of distance.

\subsection{Relation to the Metric}

The modification of distance can be absorbed into an effective metric:
\begin{equation}
  \label{eq:metric_closure}
  g_{\mu\nu} = g^{(0)}_{\mu\nu} + \alpha_G L_{\mu\nu},
\end{equation}
where $g^{(0)}_{\mu\nu}$ is the background metric derived in Vol.~1,
$L_{\mu\nu}$ is the closure tensor,
and $\alpha_G$ is a structural coupling constant.
\begin{equation}
  \boxed{\text{The metric is determined by closure density.}}
\end{equation}
This is structurally analogous to the statement that matter curves spacetime,
but no matter field has been introduced.
Geometry is determined entirely by closure structure.

\subsection{$G$ as an Effective Coupling}

The coefficient $\alpha_G$ plays the role of gravitational coupling:
\begin{equation}
  \boxed{G_{\text{eff}} \sim \alpha_G.}
\end{equation}
This constant is not fundamental.
It arises as an effective parameter of coarse-grained closure structure.
At microscopic scales, different particle closures may contribute differently;
at macroscopic scales, coarse-graining averages these contributions.
This produces an approximately universal gravitational coupling.
\begin{equation}
  \boxed{\text{The universality of } G \text{ emerges dynamically.}}
\end{equation}

\subsection{Summary}

Closure density has been identified as the source of gravitational distortion.
\begin{equation}
  \boxed{\text{Gravity = closure density.}}
\end{equation}
The modification of distance induces an effective metric,
and spacetime geometry emerges from coarse-grained closure structure.
\begin{equation}
  \boxed{\text{Spacetime emerges as coarse-grained closure geometry.}}
\end{equation}
